\begingroup\fontsize{10}{11.8}\selectfont\setlength{\abovedisplayskip}{5pt}\setlength{\belowdisplayskip}{5pt}
\subsection{Recensement des calendriers et compatibilité simultanée des biographies régionales}
Les lignes précédentes se regroupent, sans appliquer encore aucun lecteur décimal,
dans les trois biographies de blocs produites par les transitions :
\begin{align}
 B^{\rm cl}
  &=010211|012222|010211|002111|110221,\notag\\
 B^{\rm pr}
  &=201101|121221|102011|012222|102011,
 \label{act21:eq:biografias-tpk-tres-caracteres}\\
 B^{\rm au}
  &=121200|112202|121020|010210|010200.\notag
\end{align}
En effet, les paires consécutives des deux premiers tronçons sont les lignes de \(X_0,Y_0\), et celles des deux derniers sont les lignes de \(X_1,Y_1\). Ainsi, les biographies sont une autre écriture du registre de douze transitions, non une collection de préfixes conventionnels.

Pour vérifier le mot opératoire du calendrier, soit
\(c=c_1c_2c_3c_4\in\{0,1\}^4\) et, pour chaque régime
\(\chi\in\{{\rm cl},{\rm pr},{\rm au}\}\), définissons
\[
 u_0^\chi=b_0^\chi,\qquad
 u_j^\chi=u_{j-1}^\chi L_{c_j}\quad(1\le j\le4).
\]
Pour condenser le recensement, notons par
\[
 N_{\mathrm{ar}}(c):=
 \#\{(j,\chi):u_j^\chi=b_j^\chi,\ 1\le j\le4\},\qquad
 N_{\mathrm{bio}}(c):=
 \#\{\chi:(u_0^\chi,\ldots,u_4^\chi)=B^\chi\}.
\]
La multiplication exacte dans \(\FF_3\) des seize calendriers donne :
\[
\begin{array}{c|c|c}
 c&N_{\mathrm{ar}}(c)&N_{\mathrm{bio}}(c)\\ \hline
0000,0001&6&0\\
0010&9&0\\
0011&12&3\\
0100,0101,0110,0111&3&0\\
1000,1001,1010,1011&0&0\\
1100,1101,1110,1111&0&0
\end{array}
\tag{D}
\label{act21:eq:censo-dieciseis-calendarios}
\]
Le tableau épuise les \(2^4=16\) possibilités et chacune de ses entrées s'obtient
par quatre multiplications d'une ligne par l'une des deux matrices reproduites.
Par conséquent, \(0011\) est l'unique calendrier qui reconstruit
simultanément les trois biographies complètes. En notation opératorielle,
\begin{equation}
 (L_0,L_0,L_1,L_1)
 \label{act21:eq:calendario-0011-interno}
\end{equation}
ce n'est pas une prescription indépendante : c'est la seule lecture binaire compatible
avec les douze arêtes déjà produites par \(U_t\).

Le calendrier \eqref{act21:eq:calendario-0011-interno} produit quatre nouveaux
blocs de six composants. Leurs concaténations forment
\[
 w_6\longrightarrow w_{12}\longrightarrow w_{18}
 \longrightarrow w_{24}\longrightarrow w_{30}.
\]
Les sous-indices de \(w\) indiquent la longueur accumulée. La frontière \(R_{36}\) conserve la structure terminale et ouvre la relation de prolongement ; elle n'est pas renommée comme un dernier bloc périodique qui pourrait se répéter indéfiniment.

L'identité \(L=X^{-1}Y\) démontre l'unicité une fois les
transitions fixées ; elle ne remplace pas leur production par \(U_t\). Réciproquement,
\eqref{act21:eq:censo-dieciseis-calendarios} prouve dans l'article l'unicité
de l'ordre des quatre transports. La chaîne causale est ainsi close :
le recensement APP produit les régions initiales, le TRIT fixe le régime et
l'orientation, le TPK produit le registre des transitions, et ce registre
détermine \(L_0,L_1\) et \(0011\) avant toute réalisation archimédienne.


\par\endgroup
